\section{Technical Challenges, Week 03 Roadmap \& TA Consultation}
\label{sec:challenges_and_roadmap}

\subsection{Technical challenges resolved in Week 02}
During Week 02, we ran into and fixed three main engineering issues:
\begin{enumerate}[label=\textbf{\arabic*.}]
    \item \textbf{DataLoader memory stability:} When training for 100 epochs with heavy augmentations (RandAugment, Mixup, CutMix), PyTorch worker processes gradually caused memory fragmentation. We resolved this by setting \texttt{persistent\_workers=True}, allocating fixed shared memory, and caching dataset items, which kept RAM usage stable throughout long runs.
    \item \textbf{Balancing parameter counts:} To compare an isotropic model (ResMLP) with a multi-stage model (RepMLP) fairly, we adjusted layer depths and channel widths so that both models have almost the same size ($\approx 2.16\text{M} - 2.19\text{M}$ parameters, a $1.6\%$ difference).
    \item \textbf{Mixed precision training stability:} Training deep Vision MLPs from scratch in FP16 can occasionally lead to underflow during early epochs. We added dynamic loss scaling via PyTorch AMP and set gradient clipping at norm $1.0$, which prevented any numerical underflow or overflow.
\end{enumerate}

\subsection{Week 03 roadmap}
With our baseline models replicated, our plan for \textbf{Week 03} focuses on four key tasks:

\begin{itemize}[leftmargin=1.5em]
    \item \textbf{Train the proposed hybrid model:} Train and evaluate our proposed model, which adds structural re-parameterization directly into ResMLP's linear token-mixing layer. We will measure its accuracy gain on CIFAR-100 compared to the baselines.
    \item \textbf{Verify weight fusion mathematically:} Implement automated tests to confirm that folding the convolutional branches into the linear token matrix produces the exact same output:
    \begin{equation}
        \max \|\mathbf{Y}_{\text{train\_eval}} - \mathbf{Y}_{\text{deployed}}\|_{\infty} < 10^{-4}.
    \end{equation}
    \item \textbf{Profile inference speed and latency:} Measure latency (ms), throughput (images/sec), and GPU memory usage to verify that the deployed model runs at the same speed as vanilla ResMLP without extra overhead.
    \item \textbf{Scale to Tiny-ImageNet:} Run experiments on Tiny-ImageNet ($64 * 64$ resolution, 200 classes) to see how the models perform on slightly larger images.
\end{itemize}

\subsection{Questions for Teaching Assistant (TA) consultation}
In MetaFormer~\cite{yu2022metaformer}, Yu et al. showed that the general architecture (separating token mixing and channel mixing) is a key factor for vision models, demonstrating that even an extremely simple token mixer like average pooling (PoolFormer) can achieve strong performance.

However, comparing the ImageNet benchmark curves of PoolFormer and ResMLP highlights a major difference in parameter efficiency:
\begin{itemize}[leftmargin=1.5em]
    \item \textbf{PoolFormer parameter efficiency:} PoolFormer reaches $81.4\%$ Top-1 accuracy with only $\approx 31\text{M}$ parameters and $5.0\text{G}$ MACs (and reaches $82.5\%$ with $\approx 73\text{M}$ parameters and $11.6\text{G}$ MACs).
    \item \textbf{ResMLP parameter bloat:} In contrast, ResMLP requires a massive budget of $\approx 116\text{M}$ parameters (ResMLP-36) and $23.0\text{G}$ MACs just to reach $\approx 81.0\%$ Top-1 accuracy.
\end{itemize}

\textbf{Team Hypothesis:} Our team hypothesizes that this significant gap in parameter efficiency is driven by their macro-architectural design:
\begin{itemize}[leftmargin=1.5em]
    \item \textbf{Hierarchical pyramid with conv embeddings:} PoolFormer uses a 4-stage pyramid layout. Before each stage, a convolutional embedding layer downsamples the spatial grid while expanding channel capacity, which gradually shrinks image resolution and continuously injects local inductive bias.
    \item \textbf{Single-scale isotropic columnar design:} ResMLP uses a single-scale columnar structure. It embeds the image only once at the patch stem and maintains a constant token count $N$ and channel width $D$ across all layers. Without multi-scale spatial reduction and intermediate conv embeddings, ResMLP must scale up depth and width massively to learn rich visual features, resulting in excessive parameters.
\end{itemize}

Based on this observation, we would like to consult our Teaching Assistant on the following questions:
\begin{enumerate}[label=\textbf{Q\arabic*:}]
    \item \textbf{Validity of our architectural hypothesis:} From the TA's perspective, is our hypothesis sound that ResMLP's lower parameter efficiency stems primarily from its isotropic single-scale layout and the lack of multi-stage convolutional downsampling embeddings?
    \item \textbf{Direction for Week 03 modeling:} To address this limitation in our research project:
    \begin{itemize}
        \item \textit{Option A (Pyramid ResMLP):} Should we extend ResMLP into a hierarchical multi-stage architecture with convolutional downsampling layers before each stage (similar to PoolFormer or RepMLP)?
        \item \textit{Option B (Isotropic Re-parameterization):} Or should we keep the isotropic columnar structure (preserving ResMLP's simple uniform design) and focus strictly on improving data efficiency by adding structural re-parameterization directly into the linear token mixer?
    \end{itemize}
    Which direction would the TA recommend as more scientifically impactful and feasible within the AIO Topic Team scope?
\end{enumerate}
